\scanpage{83}
$K^{\bullet\bullet}$\quad Algèbre de DR à p.d. élémentaire

a) Anneau bigradué, \deleted{\unsure{comm.}} alterné pour
\inserted{``deg. ext''} à bidegr. $\geq0$

b) \deleted{\unsure{structure}}
Structure à p.d. sur l'idéal $\sum_{i+j>0}K^{ij}$

c) différentielle d'Anneau gradué (pour le grad ext) ne changeant pas le degré total, et comp. avec p.d.

d) $d:K^{1,0}=\Phi\longrightarrow K^{1,1}\unsure{=|\equiv}\Psi$ est un épi.

e) Posant $k=K^{00}$, $K^{**}$ est l'algèbre \inserted{alternée} à p.d. libre engendrée par $\Phi$ et $\Psi$
i.e.
\[
K^{**}\simeq\Gamma^*\Phi\otimes\Lambda^*\Psi
\]
[Il en résulte que $K^{ij}=0$ si $j>i$.]

f) \deleted{$\Psi$ est $k$-plat} $\Phi$ et $\Psi$ sont $k$-plats
(comme aussi $L=\operatorname{Ker}d=H^{10}(K^{\bullet\bullet})$)

Il résulte déjà de ceci que
\[
H^{*i}(K^{\bullet\bullet})=
\begin{cases}
\Gamma^*L&\text{si }i=0\\
0&\text{si }i>0
\end{cases}
\]
donc que $K^{\bullet\bullet}$ est une résolution de $\Gamma^*L$.
On va supposer le plus souvent

g) $L$ est un $k$-Module \deleted{\unsure{inversible}}\inserted{loc. lib. de rang 1}.
(Voire libre de rang 1 et trivialisé $\ldots$)
\margin{82}
